% !TeX root = ../Thesis.tex

%************************************************
\chapter{Conclusion}\label{ch:conclusion}
%************************************************
%%\glsresetall % Resets all acronyms to not used
In conclusion, this thesis investigated the performance of middleware-assisted sharded distributed databases in comparison to the baseline using a single centralized database, focusing on throughput, response time, and data distribution. By designing and implementing a middleware prototype using consistent hashing with virtual nodes, the thesis examined the key factors influencing performance and addressed the challenges and bottlenecks encountered during the implementation. The results demonstrated that, despite certain limitations, sharded distributed databases supported by a middleware achieve effective data distribution while adapting to various database node and shard configurations and can enhance scalability when designed and implemented meticulously. Future work should tackle bottlenecks and limitations of the implemented prototype of middleware-assisted sharded distributed databases, such as single-instance middleware and synchronous request processing. Incorporating multi-instance middleware architecture, asynchronous processing, support for running different queries or requests in parallel across associated databases with concurrency control mechanisms within the middleware, and adaptability to distinct types of databases would significantly enhance the middleware's functionalities. Hence, this thesis has successfully examined a middleware prototype for sharding in databases, providing valuable insights and paving the way for future works aimed at enhancing the efficiency and scalability of sharded distributed databases.